Redes Wi-Fi domésticas raramente são instaladas com base em critérios de engenharia. Em geral, o ponto de acesso é colocado onde o cabo de internet chega, e não onde a cobertura seria melhor. Como discutido em aula, o ganho de uma antena não cria potência: ele apenas redistribui a potência total em cada direção. Além disso, o padrão de radiação divulgado pelo fabricante é medido em um ambiente de teste isolado, sem paredes, móveis ou pessoas ao redor. Quando instalado em uma residência real, esse padrão deixa de valer como foi projetado: ele é recortado por paredes, refletido por metais e atenuado por corpos humanos, originando um padrão de radiação instalado que é uma hipótese, e não um dado.

Assim, a cobertura de uma rede Wi-Fi em ambientes internos não depende apenas da potência de transmissão ou do ganho nominal da antena. A distância, as paredes, as portas, os móveis, os objetos metálicos, a altura de instalação e os efeitos de multipercurso podem alterar significativamente o nível de sinal recebido em diferentes pontos da residência. Por esse motivo, uma posição visualmente central nem sempre oferece a melhor distribuição de cobertura.

Diante disso, este trabalho analisa experimentalmente a seguinte questão: \textit{a posição atual do ponto de acesso é a ideal para o ambiente estudado?} Para isso, a investigação compara duas situações. Na primeira, o roteador permanece em sua posição original. Na segunda, ele é reposicionado conforme uma hipótese construída a partir da planta do ambiente e do primeiro mapa de calor